% !TEX program = pdflatex
\documentclass[11pt,a4paper]{article}

% =========================================================================
% ACTA DE CONSTITUCIÓN DEL PROYECTO
% Proyecto: Horus Sentinel
% Gestión de Proyectos TICS (EST315) -- Semestre 2026-II -- Team Sobrevivientes
% =========================================================================

% --- PAQUETES DE CONFIGURACIÓN ---
\usepackage[utf8]{inputenc}
\usepackage[T1]{fontenc}
\usepackage[spanish,es-tabla]{babel}
% Compatibilidad con babel-spanish reciente (flechas de TikZ y signo de porcentaje)
\AtBeginDocument{\shorthandoff{<>}\spanishplainpercent}
\usepackage{geometry}
\geometry{
    top=3.6cm,
    bottom=2.3cm,
    left=2.2cm,
    right=2.2cm,
    % PARTE MODIFICABLE: headheight debe ser mayor que la altura del logo más alto (2.56cm = 73pt).
    headheight=80pt,
    headsep=15pt,
    footskip=20pt
}

\usepackage{graphicx}
\usepackage{fancyhdr}
\usepackage[table]{xcolor}
\usepackage{titlesec}
\usepackage{booktabs}
\usepackage{tabularx}
\usepackage{array}
\usepackage{amsmath}
\usepackage{enumitem}
\usepackage{microtype}
\usepackage{amssymb}
\usepackage{hyperref}

% Logos: funciona a 2 o 3 niveles de profundidad dentro de docs/latex/
\graphicspath{{../../assets/img/}{../../../assets/img/}}

% --- PALETA INSTITUCIONAL FIEI / UNA PUNO ---
\definecolor{AzulUNA}{RGB}{10, 45, 90}
\definecolor{DoradoFIEI}{RGB}{180, 130, 20}
\definecolor{GrisOscuro}{RGB}{50, 50, 50}
\definecolor{AzulClaro}{RGB}{228, 236, 246}

% --- CONFIGURACIÓN DE LISTAS ---
\setlist{topsep=2pt, itemsep=2pt, parsep=1pt}

% --- COLUMNAS DE TABLA ---
\newcolumntype{L}{>{\raggedright\arraybackslash}X}
\newcolumntype{P}[1]{>{\raggedright\arraybackslash}p{#1}}
\newcolumntype{C}[1]{>{\centering\arraybackslash}p{#1}}
\newcolumntype{Q}[1]{>{\raggedleft\arraybackslash}p{#1}}
\newcolumntype{R}{>{\raggedleft\arraybackslash}X}
\renewcommand{\arraystretch}{1.25}

% --- FORMATO DE ENLACES ---
\hypersetup{
    colorlinks=true,
    linkcolor=AzulUNA,
    urlcolor=AzulUNA,
    citecolor=AzulUNA,
    pdftitle={Acta de Constitución -- Horus Sentinel},
    pdfauthor={Team Sobrevivientes}
}

% --- FORMATO DE TÍTULOS ---
\titleformat{\section}{\color{AzulUNA}\normalfont\large\bfseries}{\thesection.}{0.7em}{}[\color{AzulUNA}\titlerule]
\titleformat{\subsection}{\color{AzulUNA}\normalfont\normalsize\bfseries}{\thesubsection.}{0.5em}{}
\titlespacing*{\section}{0pt}{9pt}{4pt}
\titlespacing*{\subsection}{0pt}{6pt}{2pt}

% --- ENCABEZADO Y PIE DE PÁGINA (fancyhdr) ---
\pagestyle{fancy}
\fancyhf{}

% ===== TAMAÑO DE LOS LOGOS (PARTE MODIFICABLE) =====
% logo_fiei.png  : height=1.45cm (lado izquierdo)
% logo_equipo.png: height=2.56cm (lado derecho; 25 % más grande porque la imagen tiene más borde blanco)
% Para cambiar el tamaño, modifique solo el valor height=... de cada \includegraphics.
% Si el logo más alto supera 2.6cm, aumente también headheight y top en la configuración de página.
\fancyhead[L]{%
    \includegraphics[height=1.45cm,keepaspectratio]{logo_fiei.png}%
}
\fancyhead[C]{%
    \scriptsize \textbf{\textcolor{AzulUNA}{UNIVERSIDAD NACIONAL DEL ALTIPLANO DE PUNO}}\\[2pt]
    \tiny \textcolor{GrisOscuro}{FACULTAD DE INGENIERÍA ESTADÍSTICA E INFORMÁTICA}\\[2pt]
    \scriptsize \textbf{\textcolor{AzulUNA}{Gestión de Proyectos TICS (EST315) \textbar\ Team Sobrevivientes}}\\[2pt]
    \footnotesize \textbf{\textcolor{DoradoFIEI}{ACTA DE CONSTITUCIÓN DEL PROYECTO}}%
}
\fancyhead[R]{%
    \includegraphics[height=2.56cm,keepaspectratio]{logo_equipo.png}%
}

\renewcommand{\headrulewidth}{0pt}
\renewcommand{\headrule}{}

\fancyfoot[L]{\scriptsize \textcolor{GrisOscuro}{Horus Sentinel}}
\fancyfoot[C]{\small \thepage}
\fancyfoot[R]{\scriptsize \textcolor{GrisOscuro}{Puno, Perú -- Semestre 2026-II}}
\renewcommand{\footrulewidth}{0.4pt}

\begin{document}

\begin{center}
    {\Large \textbf{\textcolor{AzulUNA}{ACTA DE CONSTITUCIÓN DEL PROYECTO}}}\\[3pt]
    {\normalsize \textbf{Horus Sentinel}}\\[4pt]
    \footnotesize \textbf{Fecha:} viernes 11/09/2026 (semana 3) \quad \textbar\quad \textbf{Versión:} 1.0 \quad \textbar\quad \textbf{Fase:} Inicio
\end{center}

\vspace{0.1cm}

\section{Datos del Proyecto}
\noindent\begin{tabularx}{\textwidth}{@{}P{4.2cm} L@{}}
\toprule
\textbf{Proyecto:} & Horus Sentinel \\
\textbf{Sponsor:} & Team Sobrevivientes, dueño y patrocinador del proyecto \\
\textbf{Project Manager:} & Clyde Neil Paricahua Pari \\
\textbf{Asesor:} & Ing. Renzo Apaza Cutipa, docente de Gestión de Proyectos TICS (EST315) \\
\textbf{Zona piloto:} & Puno (cercado) \\
\bottomrule
\end{tabularx}

\section{Propósito y Justificación}
En Puno no existe información clara sobre qué zonas son más peligrosas y en qué horarios, por lo que las personas eligen sus rutas y el serenazgo prioriza su patrullaje sin datos. Horus Sentinel muestra un mapa de calor de riesgo y sugiere la ruta más segura. El análisis de viabilidad muestra un costo valorizado de S/ 11,829.93 y un VAN económico de S/ +17,557.81 en 3 años, sujeto a la firma de un convenio con la municipalidad.

\section{Objetivos SMART}
{\small
\noindent\begin{tabularx}{\textwidth}{@{}C{0.7cm} L L P{2.4cm}@{}}
\toprule
\rowcolor{AzulClaro}
\textbf{N.\textsuperscript{o}} & \textbf{Objetivo} & \textbf{Indicador y meta} & \textbf{Plazo} \\
\midrule
O1 & Entregar el prototipo con mapa de calor y ruta segura sobre datos simulados del cercado. & RF-H01 a RF-H04 aceptados; ruta calculada en menos de 2 segundos. & Semana 8 (16/10) \\
O2 & Entregar el MVP con reportes ciudadanos y capas por horario. & RF-H05 a RF-H07 aceptados; 0 errores críticos. & Semana 16 (11/12) \\
O3 & Lograr que los ciudadanos del piloto confíen en la ruta sugerida. & 80\,\% o más de 10 participantes califica con 4/5 o más. & Semana 15 (04/12) \\
O4 & Asegurar el financiamiento de la operación. & Convenio firmado con la Municipalidad Provincial de Puno por S/ 500 al mes. & Semana 16 (11/12) \\
O5 & Cumplir el presupuesto. & Gasto en efectivo de S/ 3,229.30 o menos; CPI de 0.95 o más. & Semana 17 (18/12) \\
\bottomrule
\end{tabularx}
}

\section{Requisitos de Alto Nivel}
\begin{itemize}[leftmargin=*]
    \item Conjunto de incidentes simulados del cercado de Puno.
    \item Mapa de calor de riesgo por densidad (KDE).
    \item Ruta segura entre dos puntos con $A^*$ y costo $= \alpha \cdot \text{distancia} + \beta \cdot \text{riesgo}$.
    \item En el MVP: reportes ciudadanos anónimos y moderados, y capas por horario.
\end{itemize}

\section{Riesgos de Alto Nivel}
{\small
\noindent\begin{tabularx}{\textwidth}{@{}L C{1.0cm} C{1.0cm} C{1.6cm} L@{}}
\toprule
\rowcolor{AzulClaro}
\textbf{Riesgo} & \textbf{Prob.} & \textbf{Imp.} & \textbf{Severidad} & \textbf{Respuesta} \\
\midrule
El cronograma no alcanza en el semestre & 3 & 3 & 9 & Revisar la viabilidad con el cronograma detallado de la semana 4 \\
Presupuesto en efectivo superior al aporte posible del equipo & 2 & 3 & 6 & Revisar con el presupuesto detallado de la semana 4 \\
Pocos reportes ciudadanos para el MVP & 3 & 2 & 6 & Campaña de difusión en el cercado \\
Falta de experiencia en análisis geoespacial & 2 & 2 & 4 & Estudio guiado antes del primer sprint \\
\bottomrule
\end{tabularx}
}

\section{Supuestos y Restricciones}
\begin{itemize}[leftmargin=*]
    \item \textbf{Supuestos:} existe un mapa abierto actualizado de las calles del cercado; la municipalidad está dispuesta a evaluar un convenio.
    \item \textbf{Restricciones:} prioridad de CasaScore cuando ambos proyectos compitan por tiempo; solo datos simulados en la Unidad I; sin identificación de personas.
\end{itemize}

\section{Hitos}
{\small
\noindent\begin{tabularx}{\textwidth}{@{}L C{1.4cm} P{2.6cm}@{}}
\toprule
\rowcolor{AzulClaro}
\textbf{Hito} & \textbf{Semana} & \textbf{Fecha} \\
\midrule
Acta de Constitución aprobada & 3 & 11/09/2026 \\
Cronograma y presupuesto detallados & 4 & 18/09/2026 \\
Prototipo con datos simulados & 8 & 16/10/2026 \\
Piloto con ciudadanos & 15 & 04/12/2026 \\
MVP y convenio municipal & 16 & 11/12/2026 \\
\bottomrule
\end{tabularx}
}

\section{Presupuesto Resumido}
{\small
\noindent\begin{tabularx}{\textwidth}{@{}L Q{3.0cm}@{}}
\toprule
\rowcolor{AzulClaro}
\textbf{Concepto} & \textbf{S/} \\
\midrule
Estimación de costos del semestre (valorizada) & 11,829.93 \\
Contingencia por riesgos y reserva de gestión (10\,\%) & 1,833.64 \\
\textbf{Presupuesto total (BAC)} & \textbf{13,663.57} \\
Gasto en efectivo, aportado por los 4 integrantes (S/ 807.33 cada uno) & 3,229.30 \\
\bottomrule
\end{tabularx}
}

\section{Interesados Clave}
{\small
\noindent\begin{tabularx}{\textwidth}{@{}P{3.6cm} L C{1.6cm} L@{}}
\toprule
\rowcolor{AzulClaro}
\textbf{Interesado} & \textbf{Interés} & \textbf{Influencia} & \textbf{Estrategia} \\
\midrule
Team Sobrevivientes (sponsor) & Producto viable y financiado & Alta & Decisión en reunión de equipo \\
Ciudadanos y estudiantes & Desplazarse con menor riesgo & Alta & Gestionar de cerca: piloto \\
Municipalidad Provincial de Puno & Prevención del delito con datos & Alta & Gestionar de cerca: propuesta de convenio \\
Serenazgo & Priorizar el patrullaje & Media & Mantener informado \\
Docente asesor & Aplicación correcta de la gestión de proyectos & Media & Revisión en cada entrega del curso \\
\bottomrule
\end{tabularx}
}

\section{Criterios de Éxito}
\begin{enumerate}[leftmargin=*]
    \item \textbf{Principal:} satisfacción de los usuarios: 80\,\% o más de los participantes del piloto califica la ruta sugerida con 4/5 o más (O3).
    \item Prototipo en la semana 8 y MVP en la semana 16 (O1 y O2).
    \item Convenio municipal firmado (O4) y gasto en efectivo dentro de lo aprobado (O5).
\end{enumerate}

\section{Project Manager y Nivel de Autoridad}
Se designa a \textbf{Clyde Neil Paricahua Pari} como Project Manager. Puede priorizar el backlog, asignar tareas, aprobar cambios menores y usar la reserva de contingencia. Los cambios mayores, el uso de la reserva de gestión y la continuidad del proyecto los decide el equipo como sponsor.

\section{Aprobación}
\noindent En calidad de sponsor, los integrantes del equipo aprueban la presente acta; el docente asesor da su visto bueno académico:

\vspace{1.5cm}
\begin{center}
\begin{tabular}{c c c}
    \parbox{6.5cm}{\centering \rule{6.0cm}{0.6pt}\\[5pt] \textbf{Clyde Neil Paricahua Pari}\\[2pt] {\footnotesize Project Manager}}
    & \hspace{0.8cm} &
    \parbox{6.5cm}{\centering \rule{6.0cm}{0.6pt}\\[5pt] \textbf{Ruth Karina Apaza Solis}\\[2pt] {\footnotesize Analista de Planificación}}
    \\[2.0cm]
    \parbox{6.5cm}{\centering \rule{6.0cm}{0.6pt}\\[5pt] \textbf{Eliseo Tarqui Ajahuana}\\[2pt] {\footnotesize Especialista Técnico / Arquitectura}}
    & \hspace{0.8cm} &
    \parbox{6.5cm}{\centering \rule{6.0cm}{0.6pt}\\[5pt] \textbf{Viviana Mary Quisbert Quispe}\\[2pt] {\footnotesize Gestora de Comunicación, Riesgos y Documentación}}
\end{tabular}
\end{center}

\vspace{1.2cm}
\begin{center}
\parbox{8cm}{\centering \rule{7.0cm}{0.6pt}\\[5pt] \textbf{Ing. Renzo Apaza Cutipa}\\[2pt] {\footnotesize Docente asesor -- Visto bueno académico}}
\end{center}
\end{document}
